% chapters/ch16-conceptual-design-review.tex
\chapter{A Design Review Method: How to Critique a Schema}

\begin{learningobjectives}
In this chapter you will learn a structured method to review and critique a relational design.
You will be able to conduct a design review using a small set of questions, identify common failure modes, and produce actionable feedback backed by rules, constraints, and dependencies.
\end{learningobjectives}

\section{Why a review method matters}
Most schema problems are not caused by ignorance of the relational model.
They are caused by missing review.
A design can look plausible while still allowing invalid states or forcing future duplication.

A review method is valuable because it makes critique repeatable.
It also makes feedback objective: instead of “I don’t like this table”, you can say “this relation mixes facts that imply the dependency \(A \rightarrow B\) but \(A\) is not a key”.

\begin{definitionbox}
A \term{design review} is a structured evaluation of a schema and its documentation to check correctness, clarity, and maintainability.
A good review produces specific issues and specific actions.
\end{definitionbox}

\section{Inputs to a review}
A review cannot be done from a schema alone.
You need a minimal context set.

A good review package includes:
\begin{itemize}[leftmargin=*,nosep]
  \item a short domain narrative (6--12 sentences),
  \item a numbered business rule list,
  \item the logical schema in book notation,
  \item a constraint block (C1, C2, ...),
  \item 6--10 key questions the database must support.
\end{itemize}

If one of these is missing, the review should start by requesting it, because otherwise critique becomes guessing.

\section{The review flow}
Use the same flow every time.
It keeps discussions focused and prevents people from jumping to performance too early.

\begin{principlebox}
\textbf{Design review flow}\par\smallskip
\begin{enumerate}[label=\arabic*.]
  \item Domain clarity: narrative and vocabulary.
  \item Identity: keys and uniqueness.
  \item Relationships: cardinality, optionality, and association relations.
  \item Constraints: what is enforced, what is missing.
  \item Redundancy: dependencies and normalization signals.
  \item Time and state: events, history, and temporal constraints.
  \item Validation: key questions and paper test data.
\end{enumerate}
\end{principlebox}

Each step has a small set of concrete checks.

\begin{figure}[htbp]
\centering
\includegraphics[width=0.85\textwidth]{review-process.pdf}
\caption{A repeatable review flow turns critique into a controlled loop.}
\end{figure}

\section{Step 1: Domain clarity checks}
A schema cannot be correct if the domain is not precise.

Check:
\begin{itemize}[leftmargin=*,nosep]
  \item Are core terms defined (e.g., title vs copy, appointment vs visit)?
  \item Are there overloaded words (e.g., “user”, “status”, “date”) used without definition?
  \item Are the main business rules written as testable statements?
\end{itemize}

Common issue: the narrative mixes two concepts under one name.
Fix: split vocabulary first, then redesign.

\section{Step 2: Identity checks}
Identity failures are catastrophic because they break references.

Check:
\begin{itemize}[leftmargin=*,nosep]
  \item Does every entity have a primary key?
  \item Is the key stable (does it change)? Is it available at creation time?
  \item Are natural identifiers that must be unique actually constrained as unique?
  \item Are there any “keys” that are actually descriptions (e.g., name as primary key)?
\end{itemize}

Common issue: using mutable business data as a primary key.
Fix: introduce surrogate identity and constrain the natural identifier separately.

\section{Step 3: Relationship checks}
Relationship structure often reveals conceptual mistakes.

Check:
\begin{itemize}[leftmargin=*,nosep]
  \item Are 1:N relationships represented with FKs on the N side?
  \item Are N:M relationships represented with association relations?
  \item Are relationship attributes stored on the association/event relation (not on entities)?
  \item Is optionality stated and consistent with the rules?
\end{itemize}

Common issue: encoding many-to-many as a comma-separated list.
Fix: create an association relation.

\section{Step 4: Constraint checks}
Constraints make the model safe.

Check:
\begin{itemize}[leftmargin=*,nosep]
  \item Are entity integrity and referential integrity constraints stated?
  \item Are domain constraints stated (allowed sets, non-negative values, date ordering)?
  \item Are the top 3--5 business constraints stated explicitly?
  \item For each business constraint: is enforcement feasible, or is detection documented?
\end{itemize}

Common issue: the schema implies constraints in prose, but there is no constraint block.
Fix: create a numbered constraint block and map rules to constraints.

\section{Step 5: Redundancy and dependency checks}
Redundancy is often hidden.
Dependencies expose it.

Check:
\begin{itemize}[leftmargin=*,nosep]
  \item For each relation, can you state what determines each attribute?
  \item Do any non-key attributes determine other non-key attributes (transitive dependencies)?
  \item Are code/name pairs duplicated (status\_code and status\_name stored together repeatedly)?
  \item Are entity attributes stored inside association relations (2NF violations)?
\end{itemize}

Common issue: a relation mixes entity facts and event facts.
Fix: split into entity relations plus event/association relations.

\section{Step 6: Time and state checks}
Time-heavy domains break naive schemas.

Check:
\begin{itemize}[leftmargin=*,nosep]
  \item Are events modelled as entities (Loan, Appointment, Payment)?
  \item Is “active” defined precisely where it matters (active loan, active reservation)?
  \item Are temporal constraints stated (no overlap, one active at a time)?
  \item Is history preserved where needed, or is overwrite acceptable and documented?
\end{itemize}

Common issue: storing current state on an entity without event history.
Fix: introduce an event relation and decide authoritative source of truth.

\section{Step 7: Validation checks}
Validation makes critique concrete.

Check:
\begin{itemize}[leftmargin=*,nosep]
  \item Can the schema support the key questions without duplication?
  \item Can you represent two valid scenarios with small paper test data?
  \item Can you represent an invalid scenario as a clear constraint violation?
  \item Is there a rule-to-schema traceability map for the important rules?
\end{itemize}

Common issue: the schema supports the “happy path” but not exceptions (cancellations, returns, replacements).
Fix: add state/time modelling and constraints.

\section{How to write actionable review comments}
Good review comments do two things:
they identify the issue, and they propose a specific repair with justification.

A practical format:
\begin{itemize}[leftmargin=*,nosep]
  \item \textbf{Observation:} what is wrong, referencing a relation/attribute.
  \item \textbf{Risk:} what anomaly or invalid state it allows.
  \item \textbf{Fix:} what structural change to make.
  \item \textbf{Justification:} which rule/dependency/constraint motivates it.
\end{itemize}

\begin{examplebox}
\textbf{Example review comment}\par\smallskip
Observation: \rel{AppointmentRecord} stores doctor\_name and doctor\_specialty per appointment.\par
Risk: update anomaly when a doctor changes specialty; inconsistent doctor naming.\par
Fix: introduce \rel{Doctor} and reference it from \rel{Appointment}.\par
Justification: doctor\_id \(\rightarrow\) name, specialty; doctor facts should be determined by doctor identity, not by appointment.
\end{examplebox}

This tone stays technical and avoids subjective critique.

\section{A short review checklist (printable)}
Use this as a one-page checklist when reviewing schemas.

\begin{examplebox}
\textbf{Checklist}\par\smallskip
\begin{itemize}
  \item Domain: terms defined; rules testable.
  \item Keys: every entity has a PK; uniqueness constraints stated.
  \item Relationships: correct 1:N mapping; N:M via association relation; optionality explicit.
  \item Constraints: entity/ref integrity; domain checks; top business constraints.
  \item Redundancy: no mixed facts; dependencies align with keys; no transitive duplication.
  \item Time: events as entities; “active” defined; temporal constraints stated.
  \item Validation: key questions; traceability; paper test scenarios.
\end{itemize}
\end{examplebox}

\begin{chaptersummary}
A design review should be repeatable.
Start from domain clarity, then check identity, relationships, constraints, redundancy, time/state, and validation.
Use dependencies to justify structural feedback, and write review comments as observation, risk, fix, and justification.
This makes critique objective and actionable.
\end{chaptersummary}

\begin{exercisebox}
Pick one of your schemas (clinic or library) and conduct a self-review.
Answer in full sentences.
\begin{enumerate}[label=\arabic*.]
  \item Identify two domain terms that could be ambiguous and define them precisely.
  \item Identify one key decision you would challenge and propose an alternative.
  \item Identify one missing business constraint and write it as a constraint statement.
  \item Identify one redundancy risk and state the dependency that reveals it.
  \item Write three review comments using the observation/risk/fix/justification format.
\end{enumerate}
\end{exercisebox}
